\documentclass[11pt,aspectratio=169]{beamer}

\usepackage[T1]{fontenc}
\usepackage{lmodern}
\usepackage{graphicx}
\usepackage{amsmath,amssymb,mathtools}
\usepackage{booktabs}
\usepackage{array}
\usepackage{xcolor}

\setbeamercovered{transparent=5}
\setbeamertemplate{navigation symbols}{}
\setbeamertemplate{itemize items}[circle]
\setbeamertemplate{itemize subitem}{--}
\setbeamertemplate{footline}[frame number]
\setbeamerfont{framesubtitle}{size=\small}
\setbeamercolor{block title}{fg=black,bg=white!80!black}
\setbeamercolor{block body}{fg=black,bg=white!97!black}
\setbeamersize{text margin left=0.65cm,text margin right=0.65cm}
\linespread{1.12}

\definecolor{mainblue}{RGB}{38,82,122}
\definecolor{mainred}{RGB}{176,73,61}
\definecolor{softgray}{RGB}{100,100,100}
\definecolor{citecolor}{RGB}{105,105,150}
\newcommand{\blue}[1]{{\color{mainblue}#1}}
\newcommand{\red}[1]{{\color{mainred}#1}}
\newcommand{\source}[1]{\par\vspace{0.15em}{\tiny\color{citecolor}#1}}

\title{Fertility, Housing, and Population Dynamics}
\subtitle{From a stationary closure to a dated transition}
\author{Tommaso De Santo}
\date{August 2026}

\begin{document}

{
\setbeamertemplate{footline}{}
\begin{frame}[plain,c]
  \titlepage
\end{frame}
}

% ---------------------------------------------------------------------------
\begin{frame}{The household mechanism stays; the population closure changes}
\begin{columns}[T,onlytextwidth]
\begin{column}{0.48\textwidth}
\textbf{The model shown to Raquel}
\begin{itemize}
  \item Sequential birth attempts and declining fecundity
  \item Children ever born and children at home tracked separately
  \item Rent versus own, discrete house sizes, saving, income risk, and bequests
  \item Housing costs affect fertility through space needs and tenure access
\end{itemize}
\end{column}
\begin{column}{0.48\textwidth}
\textbf{The missing aggregate object}
\begin{itemize}
  \item How current births generate future households
  \item How cohort sizes evolve away from stationarity
  \item How housing clears as population changes
  \item Where the observed economy sits on that path
\end{itemize}
\end{column}
\end{columns}

\medskip
\begin{center}
Household choices $\longrightarrow$ births $\longrightarrow$ entrants
$\longrightarrow$ housing demand $\longrightarrow$ prices $\longrightarrow$ choices
\end{center}

\textbf{This is a new equilibrium closure, not a new household model.}
\end{frame}

% ---------------------------------------------------------------------------
\begin{frame}{A two-generation model isolates the feedback}
\small
\begin{columns}[T,onlytextwidth]
\begin{column}{0.52\textwidth}
Young households choose consumption, adult space, and children:
\[
u(c,h,n)=c+\alpha\log h+v\log n,
\]
\[
c+ph+\big[q+a(p+\omega)\big]n=y.
\]
Their demands are
\[
n(p)=\frac{v}{q+a(p+\omega)},
\qquad
h_Y(p)=\frac{\alpha}{p}+a n(p).
\]
\end{column}
\begin{column}{0.44\textwidth}
Old households retain housing:
\[
h_O(p)=\frac{\bar h_O}{p}.
\]
Housing supply is
\[
H^s(p)=\bar H p^\varepsilon.
\]
Demography evolves according to
\[
Y_{t+1}=\frac{n(p_t)}{\bar n}Y_t,
\qquad O_{t+1}=Y_t.
\]
\end{column}
\end{columns}

\medskip
\textbf{Economic content.} A higher housing cost raises the cost of children; fewer children reduce future household formation; a smaller population eventually lowers housing costs.

\source{Source: two-generation model in the current paper. The objects $a$, $q$, $\omega$, and $v$ summarize the full model's child-space requirement, other child costs, tenure frictions, and fertility preference.}
\end{frame}

% ---------------------------------------------------------------------------
\begin{frame}{The steady state combines replacement fertility and housing clearing}
\centering
\includegraphics[width=0.96\textwidth,height=0.68\textheight,keepaspectratio]{../output/model/fertility_population_housing_transition_note/steady_state_comparative_statics.pdf}

\vspace{-0.4em}
\small
\begin{itemize}
  \item The child decision pins down the housing cost consistent with replacement: $n(p^*)=\bar n$.
  \item Housing clearing then determines how many households can live at that cost.
\end{itemize}
\source{Blue: initial value of children. Red: permanently lower value. The arrows compare impact and steady-state outcomes; they are not the literal transition path.}
\end{frame}

% ---------------------------------------------------------------------------
\begin{frame}{Below-replacement fertility need not lower population immediately}
Let $m_t=n(p_t)/\bar n$ be the replacement ratio and $\mu_t=O_t/Y_t$ the old-to-young cohort ratio. Population can rise with $m_t<1$ whenever
\[
\mu_t<m_t<1.
\]

\begin{block}{A numerical example}
\centering
\Large
$100$ young $+$ $70$ old
$\quad\longrightarrow\quad$
$80$ young $+$ $100$ old

\normalsize
\medskip
$m_t=0.8<1$, but total households rise from $170$ to $180$.
\end{block}

Housing costs can rise as well if the aging cohort retains more housing than is released by the shrinking young cohort:
\[
(1-\mu_t)h_O(p_t)>(1-m_t)h_Y(p_t).
\]

\vspace{-0.35em}
\textbf{U.S. interpretation:} inherited cohorts can keep population and housing costs rising even after fertility falls below replacement.
\end{frame}

% ---------------------------------------------------------------------------
\begin{frame}{The full model needs a dated distribution law, not a redesign}
\small
\begin{columns}[T,onlytextwidth]
\begin{column}{0.47\textwidth}
\textbf{Keep the existing household block}
\begin{itemize}
  \item Sequential fertility and child maturation
  \item Tenure, house size, saving, and moving costs
  \item Earnings risk, retirement, survival, and bequests
\end{itemize}
\end{column}
\begin{column}{0.49\textwidth}
\textbf{Add three aggregate objects}
\begin{enumerate}
  \item A calendar-time distribution $g_t(j,x)$ over age and household states
  \item A birth-vintage queue carrying births to entry twenty years later
  \item Housing-market clearing at every date
\end{enumerate}
\end{column}
\end{columns}

\medskip
\[
B_t=\frac{b_t^{\mathrm{adjusted}}}{2.1},
\qquad
e_{t+1}=M+\rho B_{t-4},
\qquad
g_{t+1}=\mathcal Q_t g_t+\text{entrants}_{t+1}.
\]
\[
H^s(P_t)=\int h_t(j,x)\,dg_t(j,x).
\]

\textbf{The factor $2.1$ is explicit replacement accounting:} adjusted child units per future household, combining sex composition, survival to entry age, and household formation. There is no hidden child-death state.
\end{frame}

% ---------------------------------------------------------------------------
\begin{frame}{We calibrate 2007--2023 as a transition, not as a steady state}
\small
\begin{center}
\setlength{\tabcolsep}{2pt}
\begin{tabular}{>{\centering\arraybackslash}m{0.17\textwidth}c>{\centering\arraybackslash}m{0.17\textwidth}c>{\centering\arraybackslash}m{0.17\textwidth}c>{\centering\arraybackslash}m{0.17\textwidth}}
\textbf{Old benchmark} & $\longrightarrow$ & \textbf{2007 state} & $\longrightarrow$ & \textbf{2007--2023} & $\longrightarrow$ & \textbf{After 2023} \\
Completed fertility normalized to $2.1$ && Reweight only the age margin to ACS 2007 && Estimate a linear decline in the fertility preference && Hold the 2023 preference fixed and continue \\
\end{tabular}
\end{center}

\medskip
\begin{columns}[T,onlytextwidth]
\begin{column}{0.32\textwidth}
\textbf{Imposed inputs}
\begin{itemize}
  \item Census household totals
  \item ACS head-age shares
  \item 2007 start and linear path shape
\end{itemize}
\end{column}
\begin{column}{0.32\textwidth}
\textbf{Calibration}
\begin{itemize}
  \item Ten free structural parameters
  \item Twelve maintained moments
  \item Objective evaluated mainly at 2023
\end{itemize}
\end{column}
\begin{column}{0.32\textwidth}
\textbf{Historical holdouts}
\begin{itemize}
  \item Period fertility and timing
  \item Ownership and rooms
  \item House prices and rents
\end{itemize}
\end{column}
\end{columns}

\textbf{Thus 2007 is a dated initial state, not literally the old steady state, and the demographic bridge is not counted as model fit.}
\end{frame}

% ---------------------------------------------------------------------------
\begin{frame}{The 2023 fit has two important housing-size misses}
\begin{columns}[T,onlytextwidth]
\begin{column}{0.64\textwidth}
\centering
\includegraphics[width=\textwidth,height=0.66\textheight,keepaspectratio]{../output/model/e5f_no_policy_transition_report_fullhistory_roomsfix_h1_dateddid_20260817/calibrated_2023_target_fit.pdf}
\end{column}
\begin{column}{0.33\textwidth}
\footnotesize
\textbf{Calibration facts}\\[-0.15em]
\textbf{Loss: $50.742$.} Ten parameters and twelve moments.

\medskip
Two exact independent repeats.

\medskip
\textbf{Main misses}
\begin{itemize}
  \item First-birth rooms: $0.273$ versus $0.720$
  \item Mean rooms: $6.742$ versus $5.780$
\end{itemize}

They contribute $38.57$ of the total loss.

\medskip
The PSID rooms target is provisional because its prepath is not flat.
\end{column}
\end{columns}
\source{Source: final dated-transition calibration, selected candidate 7; target system \texttt{e5\_fullhistory\_roomsfix\_h1\_20260817}.}
\end{frame}

% ---------------------------------------------------------------------------
\begin{frame}{Historical validation is mixed rather than hidden}
\centering
\includegraphics[width=0.98\textwidth,height=0.69\textheight,keepaspectratio]{../output/model/e5f_no_policy_transition_report_fullhistory_roomsfix_h1_dateddid_20260817/historical_untargeted_fertility_housing_validation.pdf}

\vspace{-0.25em}
\small
\textbf{Reading the figure:} timing is close by 2023; the model misses the ownership cycle and rooms trend, and its fertility decline is too steep.

\source{National historical series are holdouts. The maintained 2023 targets use mixed samples and vintages, so this is not a unified national time-series fit.}
\end{frame}

% ---------------------------------------------------------------------------
\begin{frame}{The model does not reproduce the house-price cycle}
\centering
\includegraphics[width=0.94\textwidth,height=0.72\textheight,keepaspectratio]{../output/model/e5f_no_policy_transition_report_fullhistory_roomsfix_h1_dateddid_20260817/historical_untargeted_price_rent_validation.pdf}

\vspace{-0.2em}
\begin{itemize}
  \item The model cannot reproduce the 2007--2011 house-price collapse and later rebound.
  \item With temporary-equilibrium expectations and a fixed user-cost relation, it has no independent price-to-rent dynamics.
\end{itemize}
\textbf{The exercise should therefore be read as a demographic-housing transition, not as a complete account of the post-2007 housing cycle.}
\source{Observed house prices: real Case-Shiller. The rent CPI is shown only as a descriptive proxy, not as the model's user cost.}
\end{frame}

% ---------------------------------------------------------------------------
\begin{frame}{Demographic momentum delays the post-2023 decline}
\centering
\includegraphics[width=0.96\textwidth,height=0.69\textheight,keepaspectratio]{../output/model/e5f_no_policy_transition_report_fullhistory_roomsfix_h1_dateddid_20260817/continuation_paired_continuation_levels.pdf}

\vspace{-0.2em}
\small
Both paths peak slightly in 2027. By 2183, population is $0.138$ of 2023 in the closed benchmark and $0.359$ in the open sensitivity; asset prices are $0.386$ and $0.628$ of 2023.

\source{The open case fixes the old-state outside flow $M$ and retention rate $\rho$. The value $0.169$ normalizes the old state; it is not a constant entrant share along the path.}
\end{frame}

% ---------------------------------------------------------------------------
\begin{frame}{No positive closed steady state at the fitted 2023 preference}
\begin{columns}[T,onlytextwidth]
\begin{column}{0.61\textwidth}
\centering
\includegraphics[width=\textwidth,height=0.58\textheight,keepaspectratio]{../output/model/e5f_no_policy_transition_report_fullhistory_roomsfix_h1_dateddid_20260817/continuation_closed_stationary_renewal_schedule.pdf}
\end{column}
\begin{column}{0.36\textwidth}
\footnotesize
A closed steady state requires
\[
\frac{B(P)}{E(P)}=1.
\]
At the estimated 2023 preference,
\[
0.552\leq B(P)/E(P)\leq0.678
\]
over the declared price grid. There is no crossing.

\smallskip
Even almost free housing does not restore replacement in the fitted model.

\smallskip
\textbf{Therefore:} the blue path is a finite-horizon closed benchmark, not a verified transition between two positive steady states.
\end{column}
\end{columns}
\vspace{-0.45em}
\source{Source: 25-point grid, $P/P_0\in[10^{-4},3]$; not a global theorem.}
\end{frame}

% ---------------------------------------------------------------------------
\begin{frame}{The honest advisor summary}
\small
\begin{block}{What the update adds}
The paper keeps the household block and adds population renewal. It treats 2007--2023 as a dated transition, calibrates at 2023, and follows inherited cohorts forward.
\end{block}

\begin{itemize}
  \item \textbf{Theory:} lower fertility eventually reduces population and housing costs, but inherited cohorts can keep population and prices rising initially.
  \item \textbf{Quantitative result:} the model is credible on several 2023 margins, but weak on housing quantities and the historical price path.
  \item \textbf{Central limitation:} the fitted housing response is too weak to restore replacement, so the closed path has no positive stationary endpoint.
\end{itemize}

\smallskip
\textbf{Decision:} present the finite closed benchmark plus an explicit open sensitivity, or discipline the housing-cost-to-fertility slope and re-estimate under a closed-root requirement before making a between-steady-states claim.

\smallskip
\textbf{No policy result is promoted until that equilibrium choice is settled.}
\end{frame}

\end{document}
